\Needspace{8\baselineskip}\originalchapter{磁通守恒加热与非圆截面}\label{ch:fct}
\originalsection{加热快于电流扩散时守恒什么}
若加热时间短于磁扩散时间而长于 MHD 平衡时间, 每个阶段可近似静态, 但磁通冻结限制了通量函数. 理想轴对称演化同时保持磁面间的极向及环向磁通, 因而 q 作为两磁通导数之比保持. 这比固定总电流更强; 加热可诱导电流变化以维持通量约束. 实际装置的边界通量、电源及非感应驱动要额外指定.

课程 Lecture 11 明说两个常斜率通量函数只有两自由参数, 在加热时不能同时保持 q0 与 qa, 更不能保持整条 q(\ensuremath{\Psi}). 因而采用只固定 qa 的简化 FCT, 不称其为精确全磁通守恒演化. 设 $\nu_*=\beta_tq_a^2/\epsilon$, 则
\begin{equation}
 \nu=\nu_*(1-\nu^2),\quad
 \nu=\frac{2\nu_*}{1+\sqrt{1+4\nu_*^2}},\quad
 q_*=q_a\sqrt{1-\nu^2}.
\end{equation}
加热使 I\ensuremath{\propto}1/q* 增大, \ensuremath{\nu} 有限趋1而非在有限加热量碰到1. 这是该受约束模型的结果, 不能据此宣称真实装置没有任何平衡或稳定上限.

\originalsection{反坐标为何适合求磁面}
直接 GS 用 \ensuremath{\Psi}(R,Z), 反坐标把 R、Z 作为磁面标签 \ensuremath{\rho} 和角 \ensuremath{\mu} 的函数. 例如 $R=R_0+a[l(\rho)+\rho\cos\mu]$, $Z=a\kappa(\rho)\rho\sin\mu$; 三角形形变还可加入高谐波. Jacobian $\mathcal J=R_\rho Z_\mu-R_\mu Z_\rho$ 不得为零, 否则坐标折叠, 不再是嵌套磁面. 用反坐标, 边界和磁轴形状更直观, 但微分算符复杂, 应解析或自动求导而不能忽略 Jacobian.

有限参数近似把 GS 残差与若干角权重的积分设零. 有限维矩残差趋零只证明投影方程求解成功, 并不证明局部 GS 残差处处小. 所以还应报告强形式残差、网格加密和不同初值, 这是后面实际编程题的核心检验.
